% generated by texlua scripts/lint.lua essentials --guide (make essentials); do not edit
\definecolor{lintaccent}{RGB}{155,28,28}\definecolor{lintmute}{gray}{0.38}
\providecommand{\lintdisc}[3]{\tikz[baseline=-.6ex]{\draw[#2,line width=.45pt,fill=white] (0,0) circle (#1);\ifcase#3\relax\or\fill[#2] (0,0) -- (0,#1) arc[start angle=90,end angle=0,radius=#1] -- cycle;\or\fill[#2] (0,0) -- (0,#1) arc[start angle=90,end angle=-90,radius=#1] -- cycle;\or\fill[#2] (0,0) -- (0,#1) arc[start angle=90,end angle=-180,radius=#1] -- cycle;\or\fill[#2] (0,0) circle (#1);\fi}}
\providecommand{\lintball}[2][black]{\lintdisc{.62ex}{#1}{#2}}
\providecommand{\lintbigball}[2][black]{\lintdisc{.8ex}{#1}{#2}}
The lint places each essential on a ladder of five levels: \lintball{0} missing; \lintball{1} a placeholder, only TODO or the template's own text; \lintball{2} started, your own text but not enough of it; \lintball{3} written; \lintball{4} used where the argument needs it. An essential with no condition for its use is \lintball{4} once written. The phase in \texttt{thesis.tex} asks every essential for a level: \textsc{plan} \lintball{1}, \textsc{draft} \lintball{2}, \textsc{polish} and \textsc{final} \lintball{4}. An essential is found by its label first, then by a heading that contains one of the words given. Text is your own when it is neither a TODO line nor a line of the template, and floats, tagged lists and example boxes count toward their own essentials, not toward the text around them.

{\small\begin{xltabular}{\linewidth}{@{}>{\raggedright\arraybackslash}p{.25\linewidth}>{\raggedright\arraybackslash}X>{\raggedright\arraybackslash}p{.3\linewidth}@{}}
\toprule \textbf{Essential} & \textbf{What it is, and why} & \textbf{Written, and used} \\ \midrule \endfirsthead
\toprule \textbf{Essential} & \textbf{What it is, and why} & \textbf{Written, and used} \\ \midrule \endhead \bottomrule \endlastfoot
\multicolumn{3}{@{}l}{\textsc{Front matter}} \\
\textbf{Abstract and Kivonat}\newline{\footnotesize\color{lintmute}\texttt{\textbackslash chapter*\{Kivonat\}} and \texttt{\{Abstract\}}}\newline{\footnotesize In action: \cref{app:writing:abstract}} & The paper in one paragraph, in English and in Hungarian; the reviewer reads it first and the registration carries it.\hfill{\footnotesize\color{lintmute}(E)} & 100 own words in each.\newline Written is enough. \\
\addlinespace[.8ex]\multicolumn{3}{@{}l}{\textsc{Introduction}} \\
\textbf{Context}\newline{\footnotesize\color{lintmute}\texttt{par:context}, or a heading with \emph{context} or \emph{motivation}}\newline{\footnotesize In action: \cref{app:chapters:context}} & Why the topic matters now: a trend, a regulation or an incident, cited; relevance is scored in C.\hfill{\footnotesize\color{lintmute}(C, E)} & 40 own words.\newline Written is enough. \\
\textbf{Problem statement}\newline{\footnotesize\color{lintmute}\texttt{par:problem}, or a heading with \emph{problem}}\newline{\footnotesize In action: \cref{app:chapters:problem}} & The concrete problem, ending in a one-sentence statement; the sheet's E names the problem statement as a must.\hfill{\footnotesize\color{lintmute}(E)} & 40 own words.\newline Written is enough. \\
\textbf{Running example}\newline{\footnotesize\color{lintmute}\texttt{\textbackslash exname} and the example box}\newline{\footnotesize In action: \cref{app:chapters:problem}} & One small scenario, named once in thesis.tex (\textbackslash{}exname) and revisited in a box in every chapter, so the reader can follow each step.\hfill{\footnotesize\color{lintmute}(A)} & Named, and the introduction's box in your own words.\newline Used when an example box in every main chapter. \\
\textbf{Requirements}\newline{\footnotesize\color{lintmute}items labelled \texttt{req:}}\newline{\footnotesize In action: \cref{app:chapters:requirements}} & What a solution must do, one checkable item each; they are the rows of the gap matrix and the yardstick of the evaluation.\hfill{\footnotesize\color{lintmute}(D/1)} & Every item in your own words.\newline Used when each requirement is referenced in the related work and in the evaluation, and each derived requirement (dr:) in the evaluation. \\
\textbf{Gap, in brief}\newline{\footnotesize\color{lintmute}\texttt{par:gap}, or a heading with \emph{gap}}\newline{\footnotesize In action: \cref{app:chapters:introgap}} & That no existing approach meets the requirements together; the novelty claim, made before the full comparison.\hfill{\footnotesize\color{lintmute}(D/1)} & 40 own words.\newline Used when it points to the full comparison: the related-work chapter, the gap matrix or the positioning figure. \\
\textbf{Leading figure}\newline{\footnotesize\color{lintmute}\texttt{fig:overview}}\newline{\footnotesize In action: \cref{fig:guide:overview}} & The problem, the contributions and the headline result in one picture; the reviewer looks here first.\hfill{\footnotesize\color{lintmute}(B/1)} & Your own content, and a caption of at least a sentence.\newline Used when it is referenced from the text. \\
\textbf{Contributions}\newline{\footnotesize\color{lintmute}items labelled \texttt{c:}}\newline{\footnotesize In action: \cref{app:chapters:contributions}} & What this work claims, as checkable items; the sheet's E names stating the main contribution as a must.\hfill{\footnotesize\color{lintmute}(E)} & Every item in your own words.\newline Used when each contribution points to the chapter that delivers it and is referenced in the conclusion. \\
\addlinespace[.8ex]\multicolumn{3}{@{}l}{\textsc{Approach}} \\
\textbf{Architecture figure}\newline{\footnotesize\color{lintmute}\texttt{fig:architecture}}\newline{\footnotesize In action: \cref{fig:guide:architecture}} & The approach end to end, with what is yours set apart from what is reused.\hfill{\footnotesize\color{lintmute}(B/1)} & Your own content, and a caption of at least a sentence.\newline Used when it is referenced from the text. \\
\textbf{Threats}\newline{\footnotesize\color{lintmute}items labelled \texttt{thr:}}\newline{\footnotesize In action: \cref{app:chapters:adversary}} & Who the adversary is and what it can do; a reviewer of a security topic looks for it first.\hfill{\footnotesize\color{lintmute}(D/2)} & Every item in your own words.\newline Used when each threat is referenced in the approach, where a guarantee or an assumption meets it. \\
\textbf{Assumptions}\newline{\footnotesize\color{lintmute}items labelled \texttt{asm:}}\newline{\footnotesize In action: \cref{app:chapters:assumptions}} & Who is trusted and which properties the guarantees rest on, stated so a reviewer can check them.\hfill{\footnotesize\color{lintmute}(D/2)} & Every item in your own words.\newline Used when each assumption is referenced in the approach, by a guarantee or a proof. \\
\textbf{Guarantees}\newline{\footnotesize\color{lintmute}items labelled \texttt{g:}}\newline{\footnotesize In action: \cref{app:chapters:guarantees}} & What the approach guarantees under the assumptions, against the threats; each is what the evaluation checks.\hfill{\footnotesize\color{lintmute}(D/2)} & Every item in your own words.\newline Used when each guarantee is referenced in the evaluation, as a row of the correctness table. \\
\addlinespace[.8ex]\multicolumn{3}{@{}l}{\textsc{Evaluation}} \\
\textbf{Research questions}\newline{\footnotesize\color{lintmute}items labelled \texttt{rq:}}\newline{\footnotesize In action: \cref{app:chapters:research-questions}} & The questions the evaluation answers: correctness, scale and robustness at least.\hfill{\footnotesize\color{lintmute}(E)} & Every item in your own words.\newline Used when each research question is referenced in the evaluation, where it is answered. \\
\textbf{Setup and artifact}\newline{\footnotesize\color{lintmute}\texttt{sec:setup}, or a heading with \emph{setup}}\newline{\footnotesize In action: \cref{app:chapters:setup}} & How the measurements were made, so each can be checked, and where the archived artifact is.\hfill{\footnotesize\color{lintmute}(D/2)} & 100 own words.\newline Used when the evaluation cites the archived artifact with a real DOI. \\
\textbf{Correctness}\newline{\footnotesize\color{lintmute}\texttt{tab:correctness}}\newline{\footnotesize In action: \cref{tab:guide:correctness}} & How each guarantee and requirement was checked: proved, tested or model-checked, with the result.\hfill{\footnotesize\color{lintmute}(D/2)} & Your own content, and a caption of at least a sentence.\newline Used when it is referenced from the text, and every guarantee (g:) has a row in it. \\
\textbf{Scalability}\newline{\footnotesize\color{lintmute}\texttt{sec:scalability}, or a heading with \emph{scalab} or \emph{scaling} or \emph{performance}}\newline{\footnotesize In action: \cref{app:chapters:scalability}} & How the cost grows with the input; the review sheet scores the evaluation's depth in E.\hfill{\footnotesize\color{lintmute}(E)} & 100 own words.\newline Used when the section holds a figure or table that the text references. \\
\textbf{Robustness}\newline{\footnotesize\color{lintmute}\texttt{sec:robustness}, or a heading with \emph{robust}}\newline{\footnotesize In action: \cref{app:chapters:robustness}} & How the system behaves when its assumptions are stressed; scored apart from scalability.\hfill{\footnotesize\color{lintmute}(E)} & 100 own words.\newline Used when the section holds a figure or table that the text references. \\
\textbf{Threats to validity}\newline{\footnotesize\color{lintmute}\texttt{sec:validity}, or a heading with \emph{validity}}\newline{\footnotesize In action: \cref{app:chapters:validity}} & What could make the results wrong or not generalise, and what was done about each.\hfill{\footnotesize\color{lintmute}(D/2)} & 100 own words.\newline Written is enough. \\
\addlinespace[.8ex]\multicolumn{3}{@{}l}{\textsc{Related work}} \\
\textbf{Positioning figure}\newline{\footnotesize\color{lintmute}\texttt{fig:positioning}}\newline{\footnotesize In action: \cref{fig:guide:positioning}} & The related work on the two properties that separate it most, with this work in the cell no one else reaches.\hfill{\footnotesize\color{lintmute}(B/2)} & Your own content, and a caption of at least a sentence.\newline Used when it is referenced from the text. \\
\textbf{Gap matrix}\newline{\footnotesize\color{lintmute}\texttt{tab:related}}\newline{\footnotesize In action: \cref{tab:guide:related}} & Each requirement against the closest approaches: every other approach misses something, this work meets all.\hfill{\footnotesize\color{lintmute}(D/1)} & Your own content, and a caption of at least a sentence.\newline Used when rows are the requirements, this work's column is all met with pointers to results, and every other approach misses one. \\
\textbf{The gap}\newline{\footnotesize\color{lintmute}\texttt{sec:gap}, or a heading with \emph{gap}}\newline{\footnotesize In action: \cref{app:chapters:gap}} & The gap in one sentence, shown closed by this work's results.\hfill{\footnotesize\color{lintmute}(D/1)} & 100 own words.\newline Used when it is based on the gap matrix (references tab:related). \\
\addlinespace[.8ex]\multicolumn{3}{@{}l}{\textsc{Conclusion}} \\
\textbf{Summary with evidence}\newline{\footnotesize\color{lintmute}\texttt{sec:summary}, or a heading with \emph{summary}}\newline{\footnotesize In action: \cref{app:chapters:summary}} & Each contribution restated with the result behind it; the reviewer checks here that the promises were kept.\hfill{\footnotesize\color{lintmute}(E)} & 100 own words.\newline Used when it cites an evaluation result: the evaluation chapter, a research question, a figure, a table or a theorem. \\
\textbf{Limitations}\newline{\footnotesize\color{lintmute}items labelled \texttt{lim:}}\newline{\footnotesize In action: \cref{app:chapters:future}} & What this work does not do, stated before a reviewer finds it.\hfill{\footnotesize\color{lintmute}(E)} & Every item in your own words.\newline Written is enough. \\
\textbf{Future work}\newline{\footnotesize\color{lintmute}\texttt{sec:future}, or a heading with \emph{future}}\newline{\footnotesize In action: \cref{app:chapters:future}} & What is still open, and the next step that would lift each limitation.\hfill{\footnotesize\color{lintmute}(E)} & 100 own words.\newline Written is enough. \\
\end{xltabular}}
